\printconcepts

\exercise{A fundamental calculus technique is to use \underline{\hskip .5in} to refine approximations to get an exact answer.}{limits}

\exercise{What is the upper bound in the summation $\ds \sum_{i=7}^{14} (48i-201)$?}{14}

\exercise{This section approximates definite integrals using what geometric shape?}{Rectangles.}

\exercise{T/F: A sum using the Right Hand Rule is an example of a Riemann Sum.}{T}

\printproblems

\input{exercises/05-03-exset-01}

\input{exercises/05-03-exset-02}

\input{exercises/05-03-exset-03}

\input{exercises/05-03-exset-04}

\exercisesetinstructions{, express the limit as a definite integral.}

\exercise{$\ds \lim_{n\to\infty} \frac{\pi}{n} \sum_{i=1}^n \frac{\sin \frac{\pi i}{n}}{1+\frac{\pi i}{n}}$}{$\ds \int_0^{\pi} \frac{\sin x}{1+x}\dd x$}

\exercise{$\ds \lim_{n\to\infty} \sum_{i=1}^n \frac{3}{n} \biggl(2+\frac{3i}{n}\biggr)\sqrt{1+\biggl(2+\frac{3i}{n}\biggr)^3}$}{$\ds \int_2^5 x\sqrt{1+x^3}\dd x$}

\exercise{$\ds \lim_{n\to\infty} \frac{5}{n} \sum_{i=1}^n \biggl(5\biggl(2+\frac{5i}{n}\biggr)^3-4\biggl(2+\frac{5i}{n}\biggr)+7\biggr)$}{$\ds \int_2^7 5x^3-4x+7\dd x$}

\exercise{$\ds  \lim_{n\to\infty} \frac{2}{n} \sum_{i=1}^n \frac{1+\frac{2i}{n}}{\biggl(1+\frac{2i}{n}\biggr)^2+4}$}{$\ds \int_1^3 \frac{x}{x^2+4}\dd x$}

\exercisesetend

\exercisesetinstructions{, express the definite integral as a limit of a sum.}

\exercise{$\ds \int_2^5 4-2x\dd x$}{$\ds \lim_{n\to\infty}\left[\frac{3}{n} \sum_{i=1}^n 4-2\biggl(2+\frac{3i}{n}\biggr)\right]$}

\exercise{$\ds \int_{-2}^0 x^2+3x\dd x$}{$\ds \lim_{n\to\infty} \frac{2}{n} \sum_{i=1}^n \biggl[\biggl(-2+\frac{2i}{n}\biggr)^2+3\biggl(-2+\frac{2i}{n}\biggr)\biggr]$}

\exercise{$\ds \int_{-\pi/2}^{\pi/2} \frac{\sin^3 x}{2+\cos x}\dd x$}{$\ds \lim_{n\to\infty} \frac{\pi}{n} \sum_{i=1}^n \frac{\sin^3(-\pi/2+\pi i/n)}{2+\cos (-\pi/2+\pi i/n)}$}

\exercise{$\ds\int_0^2 e^x\dd x$}{$\ds\lim_{n\to\infty}\frac2n\sum_{i=1}^n e^{2i/n}$}

\exercisesetend

\input{exercises/05-03-exset-05}

\input{exercises/05-03-exset-06}

\exercise{Use six rectangles to approximate the area under the given graph of $f$ from $x=0$ to $x=12$, using:
\begin{enumext}
\item The Left Hand Rule,
\item The Right Hand Rule,
\item The Midpoint Rule.
\end{enumext}
\begin{tikzpicture}[alt={A parabola starting flat at (0,9) and curving downward toward a root at (12,0).}]
\begin{axis} [width=.4\textwidth,
tick label style={font=\scriptsize},axis y line=middle,axis x line=middle,
name=myplot,axis on top,%all axes={grid},
                        ymin=0,ymax=10,
                        xmin=0,xmax=12,smooth]
   \addplot [draw={\coloronefill},fill=\coloronefill,area style,domain=0:12] {9-(x/4)^2} \closedcycle;
   \addplot [smooth,thick,draw={\colorone},domain=0:12] {9-(x/4)^2};
   \draw [help lines,step=10] (axis cs:0,0) grid (axis cs:12,10);
   % steps=10 was guess and check
\end{axis}
\node [right] at (myplot.right of origin) {\scriptsize $x$};
\node [above] at (myplot.above origin) {\scriptsize $y$};
\end{tikzpicture}}{\mbox{}\\[-2\baselineskip]\parbox[t]{\linewidth}{\begin{enumext}
\item	Exact expressions will vary; $80.5$.
\item	$72.25$
\item	$62.5$
\end{enumext}}}

\exercise{A car accelerates from 0 to 40 mph in 30 seconds. The speedometer reading at each 5 second interval during this time is given in the table below. Estimate how far the car travels during this 30 second period using the velocities at:
\begin{enumext}
\item The beginning of each time interval.
\item The end of each time interval.
\end{enumext}

\begin{center}
\tagpdfsetup{table/header-columns={1}}
\begin{tabular} {|c|c|c|c|c|c|c|c|}
\hline
$t$ (sec)&0&5&10&15&20&25&30\\ \hline
$v$ (mph)&0&6&14&23&30&36&40\\ \hline
\end{tabular}
\end{center}}{\mbox{}\\[-2\baselineskip]\parbox[t]{\linewidth}{\begin{enumext}
\item	$(5\text{ s})((0+6+14+23+30+36)\text{ mph})=545\frac{\text{mi s}}{\text{hr}}\times\frac{1\text{ hr}}{3600\text{ s}}\times\frac{5280\text{ ft}}{1\text{ mi}}=799\text{ ft}$
\item	$(5\text{ s})((6+14+23+30+36+40)\text{ mph})=585\frac{\text{mi s}}{\text{hr}}\times\frac{1\text{ hr}}{3600\text{ s}}\times\frac{5280\text{ ft}}{1\text{ mi}}=858\text{ ft}$
\end{enumext}}}

\exercise{Use Theorems \ref{thm:summation} and \ref{thm:riemannSum} to justify the remaining property in \autoref{thm:defintprop}: \[\int_a^b k\cdot f(x)\dd x =k\int_a^b f(x)\dd x\]}{\mbox{}\\[-2\baselineskip]\parbox[t]{\linewidth}{\begin{align*}
\int_a^b k\cdot f(x)\dd x
&=\lim_{n\to\infty}\sum_{i=1}^n k\cdot f(c_i)\Delta x
&\iflatexml\quad\fi&\text{T\ref{thm:riemannSum}.2} \\
&=\lim_{n\to\infty}k\cdot\sum_{i=1}^n k\cdot f(c_i)\Delta x
&&\text{T\ref{thm:summation}.3} \\
&=k\cdot\lim_{n\to\infty}\sum_{i=1}^n k\cdot f(c_i)\Delta x
&&\text{T\ref{thm:limit_algebra}.4} \\
&=k\int_a^b f(x)\dd x&&\text{T\ref{thm:riemannSum}.2}
\end{align*}}}

\exercise{Use Theorems \ref{thm:summation} and \ref{thm:riemannSum} to justify the remaining property in \autoref{thm:further_def_int_props}: If $f(x)\leq M$ for all $x$ in $[a,b]$, then
\[\int_a^b f(x)\dd x\leq M(b-a).\]}{Let $f$ and $M$ be as given.
\begin{align*}
\int_a^b f(x)\dd x
&=\lim_{n\to\infty}\sum_{i=1}^n f(c_i)\Delta x
&\iflatexml\quad\fi&\text{T\ref{thm:riemannSum}.2} \\
&\le\lim_{n\to\infty}\sum_{i=1}^n M\Delta x \\
&=\int_a^b M\dd x&&\text{T\ref{thm:riemannSum}.2} \\
&=M(b-a)
\end{align*}}

\printreview

\input{exercises/05-03-exset-07}
